\section{The \texttt{Lens\_Kinoform} McXtrace Component}
A model of a specific kinoform used by the BNL

\subsection*{Identification}
\begin{itemize}
  \item \textbf{Author:} Jana Baltser and Erik Knudsen
  \item \textbf{Origin:} NBI
  \item \textbf{Date:} January 2012
\end{itemize}

\subsection*{Description}
KINOFORM. A model of a specific kinoform used by the BNL team during the APS beamtime.

\begin{verbatim}
z: [0 0.002033m]
xmin=-0.0002634
xmax=0.0002634
\end{verbatim}

the principles of the kinoform's operation are described here: http://neutrons.ornl.gov/workshops/nni\_05/presentations/min050616\_xray\_evans-lutterodt\_ken\_nni05.pdf

You may as well use a kinoform with: Lens\_parab\_Cyl(r=.5e-3,yheight=1.3e-3,xwidth=1.3e-3,d=.1e-3,N=21, material\_datafile="Be.txt")

\subsection*{Input parameters}
Parameters in \textbf{boldface} are required; the others are optional.

\begin{longtable}{p{0.22\textwidth}p{0.12\textwidth}p{0.46\textwidth}p{0.14\textwidth}}
\toprule
\textbf{Name} & \textbf{Unit} & \textbf{Description} & \textbf{Default} \\
\midrule
\endhead
datafile &  &  & "kinoform.txt" \\
material\_datafile &  &  & "Si.txt" \\
yheight & m & height of the lens. & 1e-2 \\
xwidth & m & width of the lens & 5.268e-4 \\
deltaN &  &  & 0 \\
\bottomrule
\end{longtable}

\subsection*{Links}
\begin{itemize}
  \item Component source code found in file \texttt{Lens\_Kinoform.comp}.
\end{itemize}
\IfFileExists{obsolete/Lens_Kinoform_static.tex}{\input{obsolete/Lens_Kinoform_static.tex}}{}